% Physical equations retain the supplied derivation and its local corrections.
\section{背景与物理模型}
\label{sec:background}

\subsection{磁驱柔性结构的长期动力学预测}

柔性连续体机器人通过沿结构分布的形变适应受限环境，其运动同时受到材料弹性、分布惯性及环境载荷的影响。与有限刚性关节连接的机构相比，这类系统的构型难以用少量刚体坐标完整刻画，保留哪些形变自由度，直接影响预测精度与计算规模\citep{RusTolley2015,DellaSantina2023}。对于需要反复推进状态的预测和反馈任务，物理简化、学习模型与时间积分因而共同决定了可达到的精度。

磁驱动把外部磁场与结构形变联系起来。\citet{Kim2019}利用嵌入磁性颗粒的亚毫米软连续体实现远程磁转向，\citet{KimZhao2022}从材料、建模和驱动机制总结了磁性柔性系统的设计。磁作用随构型及外加场变化，既可以在静态偏置场下改变振动特性，也可以作为受控输入与杆体交换能量。描述同一根柔性杆的不同运动工况，需要明确哪些作用已经计入能量，哪些载荷还须单独施加。

这一划分在长期振动预测中尤为重要。即使轨迹始终处于相近能级，微小频率偏差仍会随时间积累为明显的相位和构型偏差。能量接近守恒与振动时间位置准确是不同的要求，短时状态重合也不能单独说明更长时域内的误差。阻力、接触及输入进一步改变功率收支，使保守振动的分析需要与受迫响应的评价分别对应。

有限维物理模型为解释这些差别提供起点。Cosserat杆模型配合隐式时间离散和空间边值求解，已在代表性柔性与连续体算例中实现实时动力学计算\citep{Till2019}；模态或空间基函数表示则以保留部分形变自由度降低模型规模\citep{DellaSantina2023}。不同近似保留了不同的惯性、刚度及载荷耦合，因此学习与预测首先需要固定其所描述的有限维系统，随后才能比较数值推进或能量参数化的作用。

已知物理与数据残差的结合已有机器人应用。\citet{Bruder2025}将物理Koopman模型与数据残差结合，并在柔性机械臂上验证控制效果。对弱非线性振动，这一思路还引出一个比较问题：有先验网络优于从头学习，是否意味着新增残差也优于纯物理模型？残差可能描述尚未建模的非线性，也可能改变主导频率或吸收离散数据中的偏差。纯物理基线、有先验网络和无先验网络需要共享初值与时域，才能分别评价先验和新增参数的贡献。

评价还取决于观测量和参照解。杆端位移是模态位置的投影，不同分量的误差可能在这一投影中抵消；完整正则状态同时包含各模态位置与动量。相对离散训练系统的误差反映一步关系的复现程度，相对独立连续参考的误差则还包含有限步长效应。由数值轨迹训练的模型可以准确学习离散关系，却仍对连续参考积累相位偏差；某一输出误差较低，也可能与模型偏差和积分误差的抵消有关。本文同时保留这两类参照与输出，再检验其步长敏感性，使较低杆端误差与更准确的完整动力学分别得到评价。

计算预算使上述比较具有实际意义。增加参数、保留更多模态和减小步长，分别改变函数逼近、空间截断及时间推进的代价。相同参数数量也不意味着实复架构具有相同函数类或训练成本，比较因而需要同时报告精度和所用预算。本文在固定物理对象上，先用纯物理基线判断学习残差的作用，再固定已选权重改变预测步长，将模型选择联系到观测目标与预测时域。随后沿同一能量对象考察频率变化和控制插入，使长期误差的解释能够回到具体的物理与数值尺度。

\subsection{四场耦合模型及有限维表示}

考虑磁场、流体、接触与杆体力学共同作用的二维柔性杆。以参考弧长$s\in[0,L]$和时间$t$为自变量，中心线$\vec r(s,t)$描述杆体构型，局部标架描述截面取向。中心线沿弧长的变化联系轴向伸长与弯曲，材料刚度再将这些变形量联系到截面内力和内力矩。平动平衡给出惯性与分布力的关系，转动平衡则联系截面力矩及内力对形变的作用。内力和内力矩描述相邻截面的相互作用，流体、接触及磁作用则从杆外进入平衡，这一区分使内部恢复项与外部载荷分别具有明确来源。

在所用二维运动和小剪切近似下，内力分解保留伸长与弯曲贡献，略去相应剪切项，整理得到式（9）。其中，伸长项仍含构型依赖的切向刚度，弯曲项由中心线的空间变化与抗弯刚度共同决定。平面化和小剪切是这一表达的近似条件，并不将所有载荷对状态的依赖改为线性。这里的二维力学表达用于组织连续体惯性、弹性和外部作用，后文的小斜率两模态梁是进一步选定的有限维分支。

\begin{figure*}[t]
  \begin{equation}
    \EquationNineDisplay
    \tag{9}\label{eq:source9}
  \end{equation}
\end{figure*}

式（9）中，$M(s,t)$汇集杆体线质量$\rho A$与流体附加质量$\rho_fAC_m$；$C$为与相对速率$|\partial\vec r/\partial t-\vec V_f|$有关的拖曳项，其系数含$\tfrac12\rho_fDC_d$。$k_{\mathrm{tan}}$和$EI$分别描述伸长与弯曲响应，右端$\vec F$汇集来流相关项、接触、重力及磁载荷。各项投影到有限维系统后仍保留这些来源，但其能量性质须按实际作用判断。

流体拖曳以杆体相对来流的运动为依据。非零来流本身可以向杆体输入机械功，因此相对运动的耗散不意味着杆体机械能始终下降。接触载荷则由距离阈值决定是否施加，并包含法向作用与切向摩擦；运动经过接触边界时，载荷关系可能发生切换。二者在连续体方程中均为外部作用，在时间离散中却需要分别考虑其功率关系与光滑性。

磁作用同样需要按所选工况组织。可以由标量势描述的静态磁作用与弹性变形能一并进入Hamiltonian，已计入总能量的磁势不再重复加入外力势。时变输入以及一般阻力、接触作用保留为单独的广义载荷。由此，能量表示确定了保守主部，剩余载荷的位置则为随后检验非保守更新提供依据，而不由一个合并的外力符号推定所有作用均有势。

为获得广义坐标$q$的有限维系统，选取空间基函数表示中心线形变，并要求平衡残差对各保留基函数的沿杆长加权积分为零。Galerkin投影给出
\(\FrozenMath{P0126M01}\)；
允许系数依赖构型及速度时，写为
\(\FrozenMath{P0144M01}\)。
质量与拖曳项通过加权积分形成$M$和$C$，伸长与弯曲项经分部积分形成$K$的两部分，$F$保留各外载荷在基函数方向上的投影。基函数本身确定惯性及广义分量间的耦合，空间导数则进入伸长和弯曲贡献；连续体参数取常值并不自动使这些投影矩阵对角化。后文固定模态坐标与物理读出，使各预测模型接受同一个输出映射。

空间基函数的选择同时规定了近似空间及广义坐标的含义。对固定基函数，速度和加速度由时间系数的导数表示，空间变化仍由基函数承担；保留不同数目的模态会改变可表示的形变。本文在固定空间表示上检验时间收敛，因此步长加密考察的是同一有限维模型中的推进误差，与增加模态数所改变的空间误差分别解释。

边界条件还需与投影空间和端部载荷一致。固支端的位移、转角约束由选定的近似空间或自由度约束实现；分部积分留下的端部力、弯矩则进入相应边界载荷。仅在端部没有外加力和力矩时，自由端才取零剪力、零弯矩。若磁作用以端部力矩等形式施加，该项须保留在边界载荷或一致的广义力中，不能沿用未加载端部的齐次条件将其消去。本文据所选分支给定静态磁势与控制势，端部作用的表述因此需与该能量或广义力保持一致。

端部力与力矩的做功分别对应端部位移和转角变化。支承约束限制允许的形变，端部驱动则提供外部功；二者在投影中具有不同位置。因而，确定空间近似时，还应同时规定端部是否加载及其广义力表示，使内部刚度、边界条件与输入作用相互对应。

给定基函数后，杆端位移由广义坐标的固定组合获得。因而，同一杆端输出可以对应不同的模态状态，杆端误差较小也可能伴随其他位置分量或动量误差。物理模型和观测映射应一并固定，才能解释后文杆端指标与完整状态指标的不同排序，并将误差换算到同一物理长度尺度。

广义坐标与动量进一步把投影模型联系到能量表示。质量矩阵随构型变化时，正则方程包含逆质量矩阵的构型梯度，其与一般二阶机械方程的惯性对应仍须相容。本研究的变质量诊断以明确写出的Hamiltonian系统为对象；式（9）说明连续体各项的投影来源，不将一般来流、接触切换及全部耦合情形纳入该光滑诊断的结论。

学习与频率诊断采用固支自由边界、小斜率的常质量两模态磁驱梁，静态磁势保留余弦形式，阻尼、接触及时变输入取零。这一选择使式（36）数据、纯物理基线、实复网络和KAM频率诊断共享同一能量及模态映射。第\ref{sec:application}节在该显式模型上再加入标量横向磁场，评价物理端口反馈下的控制插入。输入维度由外加磁场的控制方式决定，两个模态并不因此具有彼此独立的控制通道；控制结果与学习残差的预测比较分别解释。

\subsection{相关工作与问题定位}

Hamiltonian神经网络（HNN）通过标量能量对正则状态的导数构造动力学\citep{Greydanus2019}。Deep Lagrangian Networks把Euler--Lagrange结构作为网络先验\citep{Lutter2019}，Symplectic ODE-Net将控制输入纳入Hamiltonian动力学学习，并讨论基于所学能量的控制设计\citep{Zhong2020}。前者的机器人跟踪实验与后者的受控动力学学习说明，结构先验还可组织模型输入与控制计算。其结论分别对应所设系统和驱动条件，本文将训练表示的比较与既定控制律的离散比较分开评价。这里关注同一模型中已知二次能量、剩余非线性和实复参数化的作用，以纯物理基线补充有、无先验网络之间的比较。

直接学习离散流映射提供了另一条路径。SympNets以保辛模块的复合逼近Hamiltonian系统的映射\citep{Jin2020}，结构约束直接作用于一步更新；能量网络则先给出标量函数，再通过求导和积分形成映射。两者对应的验证对象不同：连续向量场的Hamiltonian结构不保证任意积分器保辛，而映射保辛也不等于精确保持原系统能量或相位。本文因而分别检查梯度关系、一步结构残差和长期轨迹。

时间离散同时影响训练数据与预测。修正Hamiltonian和后向误差分析在相应条件下解释了辛格式的长期行为\citep{Hairer2006}。Symplectic Recurrent Neural Networks将蛙跳积分、多步训练与初态优化结合\citep{Chen2020}；\citet{Offen2022}研究适配给定辛积分器的逆修正Hamiltonian学习，说明拟合离散轨迹所得能量可能与真实连续能量有所区别。因此，模型对训练数据的准确复现应与连续动力学的恢复分别评价。

本文把一步差商与同一步中点状态配对，使监督量对应式（36）的明确离散关系。预测阶段固定已选权重改变步长，同时计算匹配步长的非线性教师和独立连续参考。这一设计将权重变化排除在加密过程之外，也保留了观察模型偏差与积分误差相互抵消的机会。不同步长下的误差变化用于判断数值推进的影响，不能直接作为新增网络参数已辨认真实非线性的依据。

生成函数为保守时间更新提供理论框架\citep[第15章]{FengQin2010}。对构型依赖质量等非分离能量，本文以保守Birkhoff表示和中点变换组织隐式格式，具体成立条件在第\ref{sec:fixed_method}节说明。阻尼与外载荷作为独立中点项加入后，完整映射的结构需要另行核验。保守流的生成函数论证、阻尼轨迹的收敛与离散功率平衡由此对应不同问题，不能相互替代。

复线性层和实虚通道激活已用于深层网络\citep{Trabelsi2018}，复值网络的模型与应用也已有系统综述\citep{Lee2022}。正则复配对$\xi=(q+ip)/\sqrt2$提供位置与动量的另一种表示，但表示形式本身不决定能量逼近能力。实复表达可以对应同一实函数，实际网络也会因隐藏层、激活和优化形成不同函数类。本文先匹配解析特征检验表示等价性，再在相同实参数预算下训练全部权重，分别观察表示与架构的作用。

两层比较均设置有、无固定二次先验的四组模型，以五项种子内配对差评价先验、表示及交互。以种子为配对单位可保留数据及优化波动，使输出均值差与重复间变化共同得到报告，而不把轨迹时间点当作独立训练重复。物理基线进一步检验新增残差的作用，杆端与完整状态指标则对应不同观测目标。实际学习取保守能量$\mathcal H=H+i0$并在实正则坐标中求梯度；一般非零虚部复生成量的表示条件属于方法推导的另一层内容，不由这组保守学习结果验证。

长期结构分析还要求将频率证据与适用定理联系。经典KAM理论涉及可积参考、Kolmogorov非退化、非共振与小扰动条件\citep{Poschel2001}，辛算法的相应结果还要求规定的解析性与格式相容性\citep[定理1]{Shang1999}。\citet[定理11]{Chen2022}在规定条件下把非零训练损失HNN的误差界联系到不变环面持续性；其中$p>4$是两自由度情形下的积分指数要求，光滑性由相应函数空间及定理假设另行规定。有限轨迹损失仍须转化为所需的统一误差界。

对本文两模态梁，二次参考的频率不随幅值变化，领先阶非退化性须从磁势高阶项考察。四次角平均给出频率变化的领先项，连续相位估计再检查这一信号在所选幅值下是否可分辨。中点离散自身也产生频移，因此需要将物理频移、数值频差及步长置于同一尺度。这里的参考初值与精确作用量、有限阶非共振筛查与完整KAM条件分别说明，使数值频率诊断保持明确的解释范围。

控制推进进一步考察输入进入一步映射的位置。给定标量磁势可产生明确的动量冲量，其对称组合采用已有分裂构造\citep{Strang1968}。本文在显式模型上使用同一物理端口反馈和输入预算，先以公共平滑外加输入比较插入精度，再观察闭环跟踪与功交换。三方案的区别在于同一势力与自由步组合的时序，并非选用三种不同反馈增益；确定性输入响应与初值块的反馈响应也分别报告，以区分插入阶数与跟踪效果。固定输入和状态反馈对应不同Jacobian，因而分别检查结构。随后将学得残差投影到同一物理端口，检验力误差改善能否转化为跟踪收益。

\subsection{研究思路与主要贡献}

本文围绕磁驱柔性杆长期预测中的误差来源，依次完成有限维能量与差分构造、实复学习比较、同模型频率诊断及控制插入分析。四步沿用共同能量对象、明确的初值与观测定义，使预测中发现的相位尺度继续在频率和控制环节接受检验，各步的比较对象和结论范围分别给出。生成函数、KAM和对称分裂结论均采用所引文献的理论；本文的贡献集中在其适用条件辨析、共享物理对象上的比较设计及相应数值发现。

\noindent（1）\emph{有限维复生成量与保守中点构造的条件衔接}：明确$\mathcal H=H+G_2-iG_1$所用的常系数、$CM^{-1}$对称及外力有势条件，区分保守生成梯度和独立中点载荷，形成矩阵左解的式（36）。以非单位标量和耦合质量矩阵检验收敛与结构，再区分阻尼中点功率平衡和连续Birkhoff二形式保持。

\noindent（2）\emph{分离表示、先验、残差与步长的双层配对比较}：以四组模型的五项配对差、物理基线、固定权重步长加密及构型依赖质量坐标检验，评价实复参数化和已知物理的作用。所测弱非线性区域中，先验提供主要精度收益；较低的离散教师误差未形成稳定的实复连续优势，细步长下杆端投影的微弱均值改善也未同时改善完整状态。

\noindent（3）\emph{同一学习算例的非线性频移与数值分辨尺度}：计算余弦磁势四次平均的领先阶作用量Hessian，以独立连续频率响应核对其幅值依赖，并比较六档步长的离散频差。中心工况离散频差约为物理频移的120倍；域内梯度误差及能量包络进一步量化学习与物理扰动，区分可计算预算与严格KAM条件。

\noindent（4）\emph{同一物理端口下三种控制插入的比较}：由标量磁势比较步内、步后和对称组合，并用原中点生成函数独立核对等价控制构造。以40条独立测试轨迹及初值匹配对照评价残差补偿，区分力误差与跟踪收益；通过连续反馈和同保持序列的两类参考，分离采样误差、积分精度与辛结构。
